\documentclass[11pt]{article}
\usepackage[margin=1in]{geometry}
\usepackage{amsmath,amssymb,amsthm}
\usepackage{booktabs}
\usepackage{array}
\usepackage{enumitem}
\usepackage[round]{natbib}
\usepackage{xcolor}
\usepackage[hidelinks]{hyperref}

\theoremstyle{plain}
\newtheorem{proposition}{Proposition}
\theoremstyle{definition}
\newtheorem{definition}{Definition}
\newtheorem{remark}{Remark}

\newcommand{\tier}{\operatorname{tier}}
\newcommand{\Lean}{\textsf{Lean}}

\title{Reliability Tiers for AI-Assisted Mathematics:\\
A Verification-Centred Framework, Failure Taxonomy, and Audit Protocol}
\author{Challenge C2 Submission\thanks{Prepared as a submission for Challenge C2
(AI for Math). The companion AI log records the literature search, every machine
verification step, and every AI-generated claim that was checked by hand.}}
\date{September 30, 2026}

\begin{document}
\maketitle

\begin{abstract}
Large language models increasingly produce mathematical arguments that read like
proofs but fail under scrutiny. The research community has responded with three
families of remedies: sampling and self-consistency, learned verifiers with
process supervision, and formal theorem proving. Practitioners still lack a way to
state what a given artefact has actually been checked for. This paper makes that
question explicit. We contribute (i) a ten-item failure taxonomy for AI-generated
mathematical claims (F1--F10) that separates statement-level defects from
inference-level, numeric, attributional, and reproducibility defects; (ii) a
six-level verification tier framework (T0--T5) in which each tier is characterised
by the guarantee it can and cannot provide, its cost class, and the defects it
retires; (iii) two propositions explaining why tiers are monotone and why no sound
and complete computable verifier for arithmetic exists, which forces any usable
framework to report residual risk rather than certify truth; and (iv) a twelve-item
audit protocol with a claim-level provenance ledger. The central practical claim is
that formal verification retires inference-soundness defects but cannot retire
statement-fidelity or provenance defects, so a reliable AI4Math pipeline must pair a
formal tier with an independent adversarial audit tier. We state the boundaries of
the framework explicitly: it is a literature-grounded analytic synthesis, and its cost
classes are order-of-magnitude estimates rather than measured timings. Where measurement
is possible we supply it: the source compiles under a pinned CI toolchain with no
undefined references and the build log is committed with the paper, and a hand-built
probe of the fidelity gate catches 14 of 16 injected statement drifts with zero false
alarms on 6 faithful controls, while exposing a blind spot for converse swaps.
\end{abstract}

\section{Introduction}

Mathematical reasoning was long treated as a domain where machine output could be
checked cheaply: a result is either proved or it is not. Recent systems have made
that intuition feel achievable. AlphaGeometry solves olympiad-level geometry
problems without human demonstration \citep{trinh2024alphageometry}, FunSearch
discovers new mathematical constructions by searching over programs
\citep{romera2024funsearch}, and neural provers interacting with the
\Lean{} proof assistant make measurable progress on formal benchmarks
\citep{zheng2021minif2f,yang2023leandojo,lin2024leanstar}.

The same literature documents a persistent reliability gap. Language models score
highly on mathematical word problems \citep{cobbe2021gsm8k,hendrycks2021math} while
remaining brittle in ways that matter: performance degrades when a problem statement
gains a semantically irrelevant clause \citep{mirzadeh2024gsmsymbolic}, intrinsic
self-correction can lower accuracy rather than raise it
\citep{huang2024selfcorrect}, models assert incorrect numerical relationships with
confidence, and generated bibliographies contain fabricated entries
\citep{ji2023hallucination}. A reader who encounters a fluent paragraph labelled
``proof'' therefore learns almost nothing about its epistemic status.

Our starting observation is that ``a proof'' denotes at least five different things
in current practice: (a) a fluent informal argument, (b) an argument that several
independent samples agree on, (c) an argument whose steps passed a learned or human
step-level review, (d) a term accepted by a proof assistant kernel, and (e) a claim
reproduced by an independent party. These are not interchangeable, yet downstream
consumers -- reviewers, students, engineering teams -- tend to read them as one.
The consequence is systematic over-trust: a claim certified at a weak level is
quoted as if it had been certified at a strong one.

This paper argues that the fix is not a better single tool but an explicit accounting
of verification level. We contribute:

\begin{enumerate}[leftmargin=*,itemsep=2pt]
\item \textbf{A failure taxonomy (F1--F10).} Ten defect classes for AI-generated
mathematical claims, organised by the layer at which the defect enters: statement,
inference, numeric, attribution, or process (Section~\ref{sec:taxonomy}).
\item \textbf{A tier framework (T0--T5).} A partial order on verification levels,
each carrying a stated guarantee, a cost class, and a coverage profile over the
taxonomy (Section~\ref{sec:framework}).
\item \textbf{Two propositions and their consequences.} Tiers are monotone by
construction (Prop.~\ref{prop:mono}); and by G\"odel's incompleteness theorem no
sound and complete computable verifier for arithmetic exists
(Prop.~\ref{prop:incomplete}), so residual risk is a structural feature of the
problem rather than an implementation defect (Section~\ref{sec:analysis}).
\item \textbf{An audit protocol.} A twelve-item checklist and a claim-level
provenance ledger that make the tier of every claim auditable after the fact
(Section~\ref{sec:protocol}).
\end{enumerate}

The framework's central operational claim is asymmetric and easy to get wrong:
formal verification retires inference-soundness defects but \emph{cannot} retire
statement-fidelity defects, because the fidelity of a formal statement to an
informal intention is itself an informal judgement. Likewise no prover kernel can
detect a fabricated citation or a contaminated benchmark. A pipeline that stops at
the formal tier therefore remains unreliable in exactly the ways that matter for
published mathematics. Section~\ref{sec:limitations} states where the framework is
weak and what evidence we did not collect.

\section{Related Work}

\subsection{Informal mathematical reasoning with language models}

Benchmark progress is the visible surface of the field. GSM8K established grade-school
word problems as a reasoning test bed \citep{cobbe2021gsm8k}, and MATH raised the
difficulty to competition level \citep{hendrycks2021math}. Chain-of-thought prompting
and self-consistency -- sampling several reasoning paths and taking a majority answer --
produced large gains and remain the default baseline
\citep{wang2023selfconsistency}. These methods optimise the probability of a correct
final answer. They provide no guarantee about the argument that produced it, and a
majority answer is evidence about a model's distribution, not about the truth of a
statement.

\subsection{Learned verifiers and process supervision}

A second family replaces sampling agreement with a learned critic. Outcome-based reward
models score final answers; process supervision scores intermediate steps, and step-level
supervision of a verifier was found to outperform outcome supervision for the purpose of
selecting correct solutions \citep{lightman2023verify}. Uesato et al.\ compare
process-based and outcome-based feedback for word problems \citep{uesato2022process}, and
Math-Shepherd removes the need for human step labels by constructing them automatically
\citep{wang2024mathshepherd}. Self-correction completes the picture only when it is driven
by an external signal: a learned corrector that consumes such a signal improves generation
\citep{welleck2022selfcorrect}, whereas intrinsic self-correction without external feedback
does not reliably improve reasoning and can degrade it \citep{huang2024selfcorrect}. Models
have partial access to their own error signal \citep{kadavath2022know}, but calibration is
imperfect, which is why we treat self-reported confidence as weak evidence and build no
tier on it.

The epistemic status of this family is statistical. A verifier that is 90\% accurate across
a distribution of steps does not certify that \emph{this} step is valid; it changes the
odds. Its failure mode is therefore correlated with the generator's: both may share the
same misconception, and neither can detect it.

\subsection{Formal theorem proving and autoformalisation}

The third family produces machine-checkable certificates. Proof search over formal
libraries began with language models trained on proof steps \citep{polu2020gptf} and was
extended by tree search over proof states \citep{lample2022hypertree}, retrieval over large
libraries \citep{yang2023leandojo}, and interleaving of informal reasoning with tactic
execution \citep{lin2024leanstar}; DeepSeek-Prover shows that large-scale synthetic data
alone moves the formal benchmark frontier \citep{xin2024deepseekprover}. Benchmarks such as
MiniF2F \citep{zheng2021minif2f} and PutnamBench \citep{tsoukalas2024putnambench} make
these claims comparable across systems, and \Lean{}~4 provides a modern kernel and tactic
language \citep{demoura2021lean4}.

Autoformalisation is both the bridge and the weak point. Wu et al.\ translate
natural-language competition problems into formal statements \citep{wu2022autoformalization},
and draft--sketch--prove decomposes an informal proof into formal sketches and steps
\citep{jiang2023draft}. In both settings the system proves a \emph{formal} sentence; whether
that sentence is the one the author meant is decided outside the kernel. At the other end of
the scale, AlphaGeometry closes geometry problems with a symbolic engine
\citep{trinh2024alphageometry}, FunSearch searches over programs and produces new
constructions \citep{romera2024funsearch}, and the Flyspeck project shows what large-scale
formalisation really costs: a decade-long effort with twenty-two listed authors to formalise
a single theorem \citep{hales2017kepler}.

\subsection{Critiques, brittleness, and provenance}

Evidence of brittleness is now systematic. Adding a semantically irrelevant clause to a word
problem can move reported accuracy substantially across model families
\citep{mirzadeh2024gsmsymbolic}; the mathematical behaviour of ChatGPT on well-defined tasks
was uneven from the start \citep{frieder2023chatgpt}; and generated text, including its
bibliographic apparatus, can be fluent and false \citep{ji2023hallucination}. Automated
research loops that write, review, and iterate without a human in the loop are now attempted
\citep{lu2024aiscientist}, which raises rather than lowers the value of an audit trail.
Architectures that pair a generator with external critics -- LLM-Modulo -- are the closest
existing relative of the tier idea \citep{kambhampati2024llmmodulo}; retrieval grounding
reduces, but does not eliminate, unsupported assertions \citep{lewis2020rag}.

What is missing from this literature is a shared currency for \emph{how much} was verified.
Papers report pass rates at whatever tier they happen to use and leave the transfer question
implicit. Our framework supplies the currency and, more importantly, states what each unit of
currency cannot buy.

\section{A Failure Taxonomy for AI-Generated Mathematical Claims}
\label{sec:taxonomy}

Before assigning levels we need a vocabulary of what can go wrong. We organise defects by the
layer at which they enter a claim: the \emph{statement} (what is being asserted), the
\emph{inference} (how it is derived), the \emph{numeric} layer (concrete quantitative
assertions), the \emph{attribution} layer (what the claim is attached to), and the
\emph{process} layer (how the result was produced and reported). Table~\ref{tab:taxonomy}
lists ten classes.

Two design choices matter. First, we separate F1 from F3 deliberately: a derivation that is
internally flawless but establishes the wrong sentence is not a proof of the intended claim,
and no formal pipeline can see the difference. Second, F7 and F8 belong to the attribution
layer rather than the inference layer, because they are falsehoods \emph{about} the work
rather than errors \emph{in} it; they survive every kernel and every verifier.

\begin{table}[t]
\centering
\small
\caption{A failure taxonomy for AI-generated mathematical claims. ``Retired at'' names the
lowest tier whose checks can in principle detect and reject the defect; in practice a check
may miss an instance, which is what the residual set records.}
\label{tab:taxonomy}
\begin{tabular}{@{}p{0.7cm}p{3.0cm}p{1.6cm}p{5.6cm}p{1.5cm}@{}}
\toprule
ID & Defect & Layer & Minimal illustration & Retired at \\
\midrule
F1 & Statement drift (misformalisation) & statement & A formal statement silently drops a
side condition, so a different theorem is proved & T5 \\
F2 & Quantifier or converse error & statement & Proving ``continuous implies
differentiable'' while reporting the converse & T4/T5 \\
F3 & Unsound inference step & inference & A step interchanges limits without justification
& T4 \\
F4 & Undischarged assumption & inference & A lemma is used outside the hypotheses under
which it was proved & T4 \\
F5 & Vacuity or triviality & inference & The hypotheses are unsatisfiable, so the theorem
has no content & T5 \\
F6 & Numeric or monotonicity slip & numeric & Asserting $\pi>22/7$, contradicted by
Archimedes' bound & T2 \\
F7 & Citation fabrication or misattribution & attribution & A bibliographic entry whose
title or identifier does not exist & T5 \\
F8 & Contamination or leakage & attribution & A benchmark result obtained on problems
present in training data & T5 \\
F9 & Over-generalisation & process & A general law asserted from a handful of examples
& T3/T5 \\
F10 & Irreproducible artefact & process & Results shipped without a pinned environment or a
runnable pipeline & T5 \\
\bottomrule
\end{tabular}
\end{table}
\section{A Tiered Verification Framework}
\label{sec:framework}

\subsection{Claims, evidence, and certification}

\begin{definition}[Claim]
A claim is a tuple $c=(S,A,P,E)$, where $S$ is the statement (informal, formal, or both),
$A$ is the set of assumptions under which $S$ is asserted, $P$ is the derivation artefact
(a prose proof, a proof script, or a program), and $E$ is the evidence bundle, that is, the
records produced by the checks that were run.
\end{definition}

\begin{definition}[Tier predicate]
For $t\in\{0,\dots,5\}$, let $V_t$ be the predicate ``$c$ passes every check of tier $t$''.
The checks are listed in Table~\ref{tab:tiers}; the check set of tier $t$ contains the check
sets of all lower tiers.
\end{definition}

\begin{definition}[Certification]
$c$ is \emph{certified at level} $t$, written $\tier(c)=t$, if $t$ is the largest index such
that $V_t(c)$ holds. Certification is always relative to a tier: we never write that a claim
is verified, only that it is certified at level $t$.
\end{definition}

\begin{definition}[Coverage and residual risk]
$\Delta(t)\subseteq\{F1,\dots,F10\}$ denotes the defect classes that the checks of levels at
most $t$ can in principle detect. The residual risk of a claim certified at level $t$ is
$R(c)=\{F1,\dots,F10\}\setminus\Delta(t)$.
\end{definition}

\begin{proposition}[Monotonicity]
\label{prop:mono}
If $V_t(c)$ holds, then $V_{t'}(c)$ holds for every $t'\le t$.
\end{proposition}
\begin{proof}
Immediate from Definition~2 by construction of the check sets. We state it because it is
routinely violated in practice: an artefact is described as ``formally verified'' on the
strength of the inference layer (F3--F4) while its statement fidelity (F1) was never reviewed.
\end{proof}

\begin{table}[t]
\centering
\small
\caption{The tier framework. Costs are order-of-magnitude classes relative to producing the
claim once: $1\times$ (no extra pass), $O(10)$, $O(10^{2})$, $O(10^{3})$ or more. The rightmost
column is the part of the guarantee that the tier does \emph{not} provide.}
\label{tab:tiers}
\begin{tabular}{@{}p{0.6cm}p{3.5cm}p{4.0cm}p{1.5cm}p{2.4cm}p{2.2cm}@{}}
\toprule
Tier & Check performed & Guarantee provided & Cost & Defects retired & Cannot retire \\
\midrule
T0 & Generation only (prose answer, no independent check) & None; fluency is not evidence
& $1\times$ & --- & F1--F10 \\
T1 & Statistical agreement: repeated sampling, voting, outcome reward models
& Confidence in the answer under the model's own distribution & $O(10)$ & partial F6
& F1--F5, F7--F10 \\
T2 & Mechanical numeric or symbolic check (rational arithmetic, computer algebra,
counterexample search) & Refutation of concrete quantitative sub-claims & $O(1)$ per check
& F6 & F1--F5, F7--F10 \\
T3 & Structured step review with explicit lemma dependencies (process reward model or human
reading) & Local plausibility of each step under stated hypotheses & $O(10^{2})$
& F3, F9 (partial) & F1, F2, F5, F7, F8, F10 \\
T4 & Formal verification: statement and proof accepted by a kernel & Derivability of the
formal sentence from the formal hypotheses & $O(10^{3})$ & F3, F4 & F1, F2 (detection),
F5, F7--F10 \\
T5 & Independent adversarial audit: fidelity review, citation and provenance check,
contamination check, reproduction from a pinned environment & Detection of statement,
attribution, and process defects & human-dependent & F1, F2, F5, F7, F8, F10
& unintended defects outside the audited set \\
\bottomrule
\end{tabular}
\end{table}

\subsection{Assigning a tier}

The procedure below is what a practitioner actually runs. It is written to fail closed: a
step that cannot be completed lowers the tier rather than being skipped.

\begin{enumerate}[leftmargin=*,itemsep=2pt]
\item \textbf{Inventory.} Enumerate every claim the artefact makes, including claims in
tables, captions, abstracts, and code comments.
\item \textbf{Statement ledger.} Record $S$ and $A$ verbatim together with the intended
informal meaning expressed in the author's own words.
\item \textbf{Numeric pass.} Identify quantitative sub-claims and test each with an
independent tool; Section~\ref{sec:analysis} gives a case in which this alone refutes a
plausible claim.
\item \textbf{Inference pass.} Attempt formalisation of each load-bearing claim. A term
accepted by a kernel certifies the inference layer. Failure of formalisation is evidence
about the derivation, not about the statement.
\item \textbf{Fidelity pass.} A party that did not write the formal statement -- human or
otherwise -- compares $S$ against the recorded intent. This gate is what retires F1 and F2.
\item \textbf{Adversarial pass.} Attempt refutation rather than confirmation: search for
counterexamples, verify every citation against a primary source, and check for
contamination.
\item \textbf{Label and disclose.} Publish the tier, the residual set $R(c)$, and the evidence
bundle.
\end{enumerate}

Steps 3 and 4 are mechanical and largely automatable. Steps 2, 5, and 6 are judgement-bearing
and must be assigned to an accountable party. The most common pipeline defect we observe is
stopping after step 4 and calling the result verified.
\section{Analysis}
\label{sec:analysis}

\subsection{Why the formal tier cannot retire statement fidelity}

\begin{remark}[Argument, not a theorem]
$\Delta(4)\cap\{F1,F2\}=\emptyset$: the checks available up to tier 4 cannot detect statement
drift or a quantifier error.
\end{remark}

We state this as an argument rather than a theorem, and give the argument explicitly. A
kernel checks a derivation relation between a formal sentence and a formal proof term.
Fidelity, by contrast, is a relation between an informal intention and a formal sentence, and
the only evidence available for it is informal. Three consequences follow. First, a successful
kernel check is consistent with having proved a drifted statement, and a failed check is
consistent with having mis-formalised a true one; the kernel is silent about which case obtains.
Second, additional prover effort does not close the gap, because the object that would have to
be checked -- the intention -- is not a formal object. Third, formalising ``the intended
meaning'' cannot discharge the obligation either: it relocates the judgement to whoever writes
and reviews the specification, without removing it. A fully formal specification of intention
would refute this remark; for open mathematical claims we know of none.

\subsection{A reproducible probe of the fidelity gate}
\label{sec:probe}

The previous subsection argues that the fidelity gate cannot be replaced by formal checks. It
does not follow that the gate is unmechanisable. To fix how far mechanisation goes, we built a
minimal checker, ran it on a hand-built probe set, and put both in the artefact
(\texttt{c2-f1-probe/}); the numbers below are reproduced in continuous integration on every
change.

The checker reads four surface features from a statement -- condition atoms drawn from a 24-word
table (\texttt{nonzero}, \texttt{positive}, \texttt{continuous}, \texttt{bounded},
\texttt{compact}, \texttt{prime}, \ldots), the quantifier class (universal or existential), the
numerals occurring in the sentence, and the count of strict inequalities -- and compares an intent
statement against its formalisation. A differing quantifier class is reported as F2, any other
difference as F1, and no difference raises no alarm. The probe set holds 22 (intent, formalisation)
pairs: 16 injected drifts across four surface classes (missing side condition 4, changed bound 3,
changed strictness 3, quantifier misplacement 4) plus 2 converse or direction swaps that are
surface-identical by construction, against 6 faithful controls. Ground-truth labels were fixed
before the checker ran.

\begin{table}[t]
\centering
\small
\caption{Probe of the mechanical fidelity check (\texttt{python run-probe.py}). Proportions
describe this checker on this hand-built probe set, not a population.}
\label{tab:probe}
\begin{tabular}{@{}lr@{}}
\toprule
Quantity & Value \\
\midrule
pairs / drifted / faithful controls & 22 / 16 / 6 \\
injected drifts detected & 14 \\
coverage on injected drifts & 14/16 = 0.875 \\
coverage, F1 class (surface changes) & 10/10 = 1.000 \\
coverage, F2 class & 4/6 = 0.667 \\
false alarms on faithful controls & 0/6 = 0.000 \\
precision when an alarm is raised & 14/14 = 1.000 \\
defect class correct when flagged & 14/14 = 1.000 \\
\bottomrule
\end{tabular}
\end{table}

Two findings matter for the framework. First, the whole surface-visible class of F1 is
mechanisable at negligible precision cost: all ten cases of missing side conditions, changed
bounds, and changed strictness were flagged, with zero false alarms on the controls, so the check
can run in a pipeline without adding review noise. Second, the F2 class split. Quantifier
misplacement was caught in all four cases, but both converse and direction swaps were missed, and
the reason is structural rather than an implementation gap: ``differentiable implies continuous''
and ``continuous implies differentiable'' differ in no atom, no numeral, no quantifier class, and
no inequality count. This is the empirical form of the argument above -- no check over the sentence
alone can separate a claim from its converse -- which is why the gate must still be signed by a
party that did not write the formalisation.

The same standard applies to the artefact itself. \texttt{paper.tex} is compiled in continuous
integration on every change to the source or the bibliography, and the PDF, the full \LaTeX{} log,
and a machine-generated evidence file are committed under \texttt{build/}, so the claim that the
paper compiles can be checked without trusting the authors.

\subsection{Three worked examples}

\textbf{Example 1 (F1, invisible to the kernel).} The intended claim is that for every odd
prime $p$ the integer $p^{2}-1$ is divisible by $8$. A pipeline formalises the side condition
away:
\[
\forall p\in\mathbb{N},\qquad \mathrm{Prime}(p)\;\Longrightarrow\;8 \mid (p^{2}-1).
\]
This sentence is false: $p=2$ gives $p^{2}-1=3$ and $8\nmid 3$. The prover fails, and the
failure carries no information about the intended claim. If instead the pipeline silently
strengthens the hypothesis until something provable appears, the artefact ships a theorem no
reader asked for. Both paths pass through the kernel without touching the defect. The intended
claim is true, and its one-line reason is worth recording precisely because it shows how much
content a drift can hide: for odd $p$, the integers $p-1$ and $p+1$ are consecutive even
numbers, so one is divisible by $4$ and the other by $2$, hence $8\mid(p-1)(p+1)=p^{2}-1$.

\textbf{Example 2 (F6, refuted by a numeric pass).} A generated discussion asserts
$\pi>22/7$. Archimedes' bounds give $223/71<\pi<22/7$, so the assertion is false; the two
values are $22/7=3.142857\ldots$ and $\pi=3.141592\ldots$. A single rational-arithmetic check
refutes the claim, while the familiarity of $22/7$ as an approximation of $\pi$ makes the error
easy to miss by reading. The refutation is self-contained: no external dataset is required.

\textbf{Example 3 (F3, refuted by formalisation).} A prose proof concludes that every continuous
real function is differentiable, via a step that assumes the difference quotient has a limit at
every point. Formalisation does not merely disagree with the proof; it fails, because the
assumed limit does not exist for $x\mapsto|x|$ at $0$. Here the kernel earns its cost -- which is
exactly why its silence about Examples 1 and 2 deserves to be stated rather than assumed.

\subsection{Why residual risk is structural}

\begin{proposition}[No sound and complete computable verifier for arithmetic]
\label{prop:incomplete}
Let $\mathrm{Tr}$ be the set of true sentences of first-order arithmetic. There is no computable
procedure $M$ whose accepted set $A(M)$ satisfies both $A(M)\subseteq\mathrm{Tr}$ (soundness) and
$A(M)=\mathrm{Tr}$ (completeness).
\end{proposition}
\begin{proof}[Proof sketch]
If $A(M)=\mathrm{Tr}$, then $\mathrm{Tr}$ would be recursively enumerable, since $M$ enumerates
it by dovetailing over all inputs. By G\"odel's first incompleteness theorem, every consistent,
recursively axiomatised theory extending Robinson arithmetic is incomplete, and the standard
argument shows that $\mathrm{Tr}$ is not recursively enumerable. The two claims contradict each
other, so no such $M$ exists.
\end{proof}

The practical reading is not defeatism. It says that ``certified'' must always be read relative
to a check set, and that the honest output of a verification pipeline is a tier plus a residual
set, never the word \emph{true}. A writing environment that refuses to render a quoted claim
without its tier label would have flagged every example in this section.

\subsection{Coverage versus cost}

Table~\ref{tab:coverage} maps the method families of Section~\ref{sec:framework} onto the
framework. Three patterns are visible. Sampling and voting buy little: they reduce the variance of
the answer, not the probability that the argument is wrong. Process supervision buys the inference
layer but shares priors with the generator, so its errors correlate with the errors it is meant to
catch. Formal proving buys the inference layer outright and is the only family that does so, at a
cost that scales with the distance between the informal statement and the available library -- a
decade of work by twenty-two listed authors for one theorem in the extreme case
\citep{hales2017kepler}. Nothing in the table retires F1, F5, F7, F8, or F10 except an independent
audit, which is why the audit sits at the top of the partial order and not at the bottom.

\begin{table}[t]
\centering
\small
\caption{Method families mapped to the tier framework. ``Partial'' marks coverage that is
statistical or restricted to a sub-class of claims; the residual set is the taxonomy minus the
retired column.}
\label{tab:coverage}
\begin{tabular}{@{}p{4.1cm}p{1.1cm}p{3.2cm}p{4.4cm}@{}}
\toprule
Method family & Tier & Defects retired & Residual risk \\
\midrule
Majority vote and self-consistency \citep{wang2023selfconsistency} & T1 & F6 partial
& F1--F5, F7--F10 \\
Outcome reward model \citep{cobbe2021gsm8k} & T1 & F6 partial & F1--F5, F7--F10 \\
Process reward model, step supervision \citep{lightman2023verify,uesato2022process,wang2024mathshepherd}
& T3 & F3, F9 partial & F1, F2, F5, F7, F8, F10 \\
Self-correction with an external corrector \citep{welleck2022selfcorrect} & T2/T3
& F3 partial, F6 partial & F1, F2, F5, F7, F8, F10 \\
Intrinsic self-correction without external feedback \citep{huang2024selfcorrect} & T0/T1
& none reliable & all \\
Retrieval grounding \citep{lewis2020rag} & T2 & F7 partial & F1, F3--F6, F8--F10 \\
Autoformalisation plus kernel check \citep{wu2022autoformalization,jiang2023draft,yang2023leandojo}
& T4 & F3, F4 & F1, F2 detection, F5, F7--F10 \\
Generator plus external critics (LLM-Modulo) \citep{kambhampati2024llmmodulo} & T3$+$T5
& F3, F4, F7, F9 & F1, F2 detection, F5, F8, F10 \\
Independent reproduction and audit & T5 & F1, F2, F5, F7, F8, F10
& defects outside the audited set \\
\bottomrule
\end{tabular}
\end{table}
\section{An Audit Protocol}
\label{sec:protocol}

The framework becomes usable only when it leaves the paper, so we state a twelve-item protocol
ordered so that mechanical checks come first and judgement-bearing gates come after.
Table~\ref{tab:protocol} lists the items together with the defects each one retires. Item~3 is
the only item that retires F1, and item~8 is the only item that retires F7.

\begin{table}[t]
\centering
\small
\caption{The twelve-item audit protocol. ``Kind'' distinguishes mechanical checks from
judgement-bearing gates; the latter are the items that must be signed by an accountable party.}
\label{tab:protocol}
\begin{tabular}{@{}p{0.6cm}p{7.5cm}p{1.9cm}p{3.0cm}@{}}
\toprule
\# & Item & Kind & Retires \\
\midrule
1 & Claim inventory: every claim in the artefact, including captions, abstracts, and code
comments, is enumerated & mechanical & --- \\
2 & Statement ledger: statement and assumptions recorded verbatim with the intended informal
meaning in the author's words & judgement & prerequisite for F1 \\
3 & Fidelity review by a party that did not write the formal statement & judgement & F1, F2 \\
4 & Assumption discharge list: every imported lemma recorded with the hypotheses under which it
holds & mechanical & F4 \\
5 & Kernel record: prover name, version, commit hash, and the exact statement file & mechanical
& F3, F4 \\
6 & Numeric spot checks with an independent tool, including at least one bound or inequality
& mechanical & F6 \\
7 & Citation verification against a primary identifier (DOI or archive record), one entry at a
time, with the checked identifier recorded & mechanical & F7 partial \\
8 & Contamination check: overlap between the evaluated set and available training data, with the
result disclosed either way & judgement & F7, F8 \\
9 & Evidence-to-claim linkage: each claim tagged with its tier and residual set & mechanical
& --- \\
10 & Environment pins for reproduction: versions, seeds, and a runnable entry point & mechanical
& F10 \\
11 & Adversarial pass: an explicit attempt to refute each load-bearing claim rather than to
confirm it & judgement & F5, F9 \\
12 & Disclosure of AI involvement and of failed attempts, including abandoned branches
& judgement & F9, F10 \\
\bottomrule
\end{tabular}
\end{table}

The protocol is backed by a ledger rather than by good intentions. For each claim we require one
row recording the claim identifier, the statement and assumptions verbatim, the tier attained,
the residual set, the pointer to the evidence bundle, and the name of the party who signed the
fidelity gate (Appendix~\ref{app:ledger}). Two properties do the real work. Every claim appears
exactly once, so a claim cannot quietly migrate to a stronger tier between drafts; and
claims without a row are visible, so absence of evidence is distinguishable from evidence of
absence.

\section{Limitations}
\label{sec:limitations}

We state the boundaries of what this paper establishes.

\begin{enumerate}[leftmargin=*,itemsep=2pt]
\item \textbf{One probe, not a controlled measurement.} The framework is a literature-grounded
analytic synthesis. The only measurement in this paper is the fidelity probe of
Section~\ref{sec:probe}: 22 hand-built pairs, not a random sample, so its proportions describe
this checker on this probe set and must not be read as defect rates. We report no inter-annotator
agreement and no measured timings; the cost classes in Table~\ref{tab:tiers} remain
order-of-magnitude estimates, and a reader who needs rates must collect them.
\item \textbf{The fidelity gate is the weakest joint.} Item~3 of the protocol carries the largest
share of the residual risk. Section~\ref{sec:probe} shows the gate is partly mechanisable: the
surface-visible part of F1 fell to a 24-word checker at zero false-alarm cost, so a machine check
narrows the gate rather than closing it. We still did not measure how often independent reviewers
agree on fidelity judgements, and a disagreement-prone gate would undermine the tier it is meant
to support.
\item \textbf{The taxonomy is not proven exhaustive.} F1--F10 are grounded in the literature we
cite and in the examples we worked, but new model families and new tool interfaces will produce
defect classes we have not named. The framework tolerates this by construction -- an unnamed defect
simply belongs to the residual set -- but the taxonomy itself is not closed.
\item \textbf{Proposition~\ref{prop:incomplete} is limitative, not constructive.} It explains why
residual risk cannot be driven to zero. It does not tell a practitioner which checks to run, and it
says nothing about non-arithmetic domains such as geometry, where diagrammatic evidence plays a
role this taxonomy does not model.
\item \textbf{Our own citation checking is bounded.} Every bibliographic entry in this paper was
checked against a primary identifier at write time, but we did not re-derive the results we
attribute to those works.
\item \textbf{Automated review is not a substitute for the audit.} Nothing here licenses replacing
human sign-off of the fidelity gate with another language model: a model auditing a model shares
the failure distribution that motivated the audit.
\item \textbf{Scope of the examples.} The three worked examples are didactic and arithmetic or
analytic in character. They establish that the defects exist and are detectable at the stated
tiers; they are not a sample from which defect frequencies can be estimated.
\end{enumerate}

\section{Conclusion}

The reliability problem in AI-assisted mathematics is usually framed as a capability problem:
models are not yet good enough at proof. That framing hides a second problem which better models
do not solve. Verification is layered, the layers are not interchangeable, and the strongest
mechanised layer on offer -- formal proof -- is silent about whether the sentence it certified is
the sentence that was meant. We proposed a ten-item failure taxonomy, a six-level tier framework
with an explicit coverage profile, two propositions that fix the tier order and explain why
residual risk is structural, a twelve-item audit protocol backed by a claim-level ledger, and a
reproducible probe that fixes how far the fidelity gate can be mechanised.

The framework reduces to three instructions. Report the tier, never the adjective. Pair the formal
tier with an adversarial audit, because they retire disjoint defect classes. Treat ``certified''
as relative to a check set, and publish the residual set wherever a claim is quoted outside its own
paper.

Two lines of work follow. The empirical one is to grow the fidelity probe of
Section~\ref{sec:probe} from a hand-built set into a sampled one and to measure inter-reviewer
agreement for the gate, the framework's weakest joint; the probe already shows where that
measurement must concentrate, namely on converse and direction errors, which no check over the
sentence alone can reach. The constructive one is to build tooling that refuses to render a claim
without a tier label, turning the ledger into a property of the writing environment rather than a
task performed afterwards.

\bibliographystyle{plainnat}
\bibliography{references}

\appendix
\section{Claim-level provenance ledger}
\label{app:ledger}

\begin{table}[h]
\centering
\small
\caption{Ledger template. One row per claim; the fidelity column is signed by a party other than
the author of the statement.}
\begin{tabular}{@{}p{1.3cm}p{3.6cm}p{2.6cm}p{1.4cm}p{2.4cm}p{2.0cm}@{}}
\toprule
Claim ID & Statement and assumptions (verbatim) & Tier & Residual set & Evidence pointer
& Fidelity signed by \\
\midrule
C-01 & \textit{(statement)} & T4 & F1, F5, F7, F8, F10 & \textit{(path or hash)} & \textit{(name)} \\
C-02 & & & & & \\
C-03 & & & & & \\
\bottomrule
\end{tabular}
\end{table}

\end{document}
